\section{Introduction}
\label{sec:intro}

Ask any major large language model whether a seahorse emoji exists, and it will confidently tell you that it does---often providing a specific Unicode codepoint, a release version, and even a visual description.
None of this is true.
There is no seahorse emoji in Unicode, and there never has been \cite{unicode2024emoji}.
Yet this ``false memory'' is remarkably persistent across model families and sizes \cite{vogel2026seahorse}, raising a fundamental question: {\bf why do LLMs confidently assert the existence of things that do not exist?}

This matters because LLMs are increasingly used as primary knowledge sources.
When a model fabricates not just an answer but detailed supporting evidence---a specific codepoint, a Unicode version number, even a visual rendering of the wrong emoji---the result is qualitatively different from typical hallucination.
It resembles the \emph{confabulation} studied in cognitive psychology, where patients produce detailed but entirely fabricated accounts that they genuinely believe to be true \cite{schacter2001seven}.
Understanding when and why LLMs confabulate is essential for building trustworthy AI systems.

Prior work has established that LLMs exhibit analogues of human memory phenomena, including Deese-Roediger-McDermott (DRM) false memories where semantically related lures are falsely recalled \cite{cao2025memory}, and collective ``Mandela effects'' that emerge and persist in multi-agent systems \cite{xu2026mandela}.
Theoretical work has shown that hallucination is statistically inevitable for general-purpose language models \cite{kalai2025hallucinate, xu2024inevitable}.
However, no study has systematically tested the specific mechanism we hypothesize: that LLMs perform \emph{category completion}, falsely asserting the existence of plausible members within well-populated categories because the semantic neighborhood is dense with real examples.

We test this hypothesis through controlled experiments probing two commercial LLMs across 162~items in 7~categories.
Our results reveal a more nuanced picture than simple category completion.
Plausible-but-nonexistent items do show elevated false positive rates (12--14\% vs.\ 0--5\% for implausible items), but category density is \emph{negatively} correlated with error rates.
Instead, item-level semantic plausibility and prompt framing emerge as the primary drivers: adversarial prompts that presuppose existence amplify false positive rates from 30\% to 80\%.
The effect is also domain-specific---emojis are highly vulnerable while structured factual domains are robust.

We make the following contributions:
\begin{itemize}[leftmargin=*,itemsep=0pt,topsep=0pt]
    \item We introduce the concept of \emph{category-completion false memories} in LLMs and design a systematic experimental framework to test it across multiple categories, plausibility levels, and prompt framings.
    \item We demonstrate that semantic plausibility and adversarial framing, not category density, are the primary drivers of LLM false memories---overturning the intuitive category-density hypothesis.
    \item We show that false memory vulnerability is domain-specific: emojis constitute a uniquely vulnerable knowledge domain compared to structured factual knowledge like country capitals and chemical elements.
    \item We provide a signal detection theory analysis revealing that models achieve reasonable discrimination ($d' = 1.8$--$2.6$) but exhibit qualitatively different failure modes under adversarial framing.
\end{itemize}
